\documentclass[../../../include/open-logic-section]{subfiles}

\begin{document}

\olfileid{sth}{cardinals}{cardsasords}
\olsection{ક્રમવાચક સંખ્યાઓ તરીકે ગણાંકો}

હકીકતમાં આપણો ગણાંકોનો સિદ્ધાંત ક્રમવાચક સંખ્યાઓના
સિદ્ધાંતનો જ (નિઃસંકોચ) ઉપયોગ કરશે. એટલે કે ગણાંકોને
કેટલીક ચોક્કસ ક્રમવાચક સંખ્યાઓ તરીકે જ વ્યાખ્યાયિત કરીશું.
ખાસ કરીને નીચેની વ્યાખ્યા આપીશું:

\begin{defn}\ollabel{defcardinalasordinal}
જો $A$ને સુક્રમિત કરી શકાય, તો $\card{A}$ એવી લઘુતમ
ક્રમવાચક સંખ્યા $\gamma$ છે કે $\cardeq{A}{\gamma}$.
કોઈ પણ ક્રમવાચક સંખ્યા $\gamma$ માટે કહીએ કે $\gamma$
એ \emph{ગણાંક} છે, જો અને માત્ર જો
$\gamma = \card{\gamma}$.
\end{defn}

હમણાં આપણે ``$A$ને સુક્રમિત કરી શકાય'' એવો શબ્દપ્રયોગ
કર્યો. ગણિતમાં લગભગ હંમેશની જેમ અહીં ``શકાય'' કહેવું
એ અસ્તિત્વના વિધાનને ઢીલી રીતે કહેવા બરાબર છે:
``$A$ને સુક્રમિત કરી શકાય'' એટલે ``$A$ને સુક્રમિત
કરે એવો કોઈ સંબંધ છે''.

પરંતુ \olref{defcardinalasordinal}માં એક મુશ્કેલી છે.
દરેક ગણનું કદ હોય, એટલે દરેક~$A$ માટે $\card{A}$
અસ્તિત્વમાં હોય, એમ આપણે ઇચ્છીએ છીએ. જોકે હમણાંની
વ્યાખ્યા એક શરતથી શરૂ થાય છે: ``\emph{જો} $A$ને
સુક્રમિત કરી શકાય\ldots''. કોઈ એવો ગણ $A$ હોય
જેને સુક્રમિત ન કરી શકાય, તો આપણી વ્યાખ્યા
$\card{A}$ નામની કોઈ વસ્તુ નક્કી જ નહીં કરે.

તેથી \olref{defcardinalasordinal} વાપરવા માટે
દરેક ગણને સુક્રમિત કરી શકાય તેની ખાતરી જોઈએ.
પરંતુ આ ખાતરી~$\ZF$માં મળતી નથી. તેથી
\olref{defcardinalasordinal}
વાપરવી હોય, તો નીચેની જેમ નવું સ્વયંસિદ્ધિ ઉમેર્યા
વિના રસ્તો નથી:
\begin{axiom}[સુક્રમિતતા]
દરેક ગણને સુક્રમિત કરી શકાય.
\end{axiom}
સુક્રમિતતાનું સ્વયંસિદ્ધિ સ્વીકાર્ય છે કે નહીં તેની
ચર્ચા \olref[choice][]{chap}માં કરીશું. હમણાંથી તો
તેનો ઉપયોગ કરીશું. તેનો ઉપયોગ કરતાં
\olref{defcardinalasordinal} મુજબ વ્યાખ્યાયિત
ગણાંકો અસ્તિત્વમાં છે અને ઇચ્છિત રીતે વર્તે છે તે
સાબિત કરવું સહેલું છે:

\begin{lem}\ollabel{lem:CardinalsExist}
દરેક ગણ $A$ માટે:
\begin{enumerate}
	\item\ollabel{cardaexists} $\card{A}$ અસ્તિત્વમાં છે અને અનન્ય છે;
	\item\ollabel{cardaapprox}  $\cardeq{\card{A}}{A}$;
	\item\ollabel{cardaidem}  $\card{A}$ ગણાંક છે, એટલે કે
	$\card{A} = \card{\card{A}}$;
\end{enumerate}
\end{lem}

\begin{proof}
$A$ નક્કી કરો. સુક્રમિતતાથી $\tuple{A, R}$ એવો
સુક્રમ મળે છે. \olref[ordinals][ordtype]{thmOrdinalRepresentation}
મુજબ $\tuple{A, R}$ અનન્ય ક્રમવાચક સંખ્યા $\beta$
સાથે ક્રમ-એકરૂપ છે. તેથી $\cardeq{A}{\beta}$.
અનંતાતીત અનુમાનથી એવી અનન્ય લઘુતમ ક્રમવાચક સંખ્યા
$\gamma$ મળે છે કે $\cardeq{A}{\gamma}$.
આથી $\card{A} = \gamma$, જેનાથી \olref{cardaexists}
અને \olref{cardaapprox} સાબિત થાય છે.
\olref{cardaidem} માટે નોંધો કે $\delta \in \gamma$
હોય, તો $\gamma$ની પસંદગીથી $\cardless{\delta}{A}$;
અને ગણસામ્ય સામ્ય સંબંધ હોવાથી
(\olref[sfr][siz][equ]{equinumerosityisequi})
$\cardless{\delta}{\gamma}$ પણ મળે છે.
તેથી $\gamma = \card{\gamma}$.
%So, for reductio, suppose that there is some ordinal $\delta \in \gamma$ such that $\cardeq{\gamma}{\delta}$. Then, $\cardeq{\cardeq{A}{\gamma}}{\delta}$ so that $\cardeq{A}{\delta}$ by \olref[sfr][set][equ]{equinumerosityisequi}, which contradicts the choice of $\gamma$ as the least ordinal such that $\cardeq{A}{\gamma}$.
\end{proof}

આગળનું પરિણામ કૅન્ટરનો સિદ્ધાંત અને તે ઉપરાંત
બીજી બાબતોની પણ ખાતરી આપે છે. (નોંધો કે ગણાંકોનો
ક્રમ તેમને ક્રમવાચક સંખ્યાઓમાંથી મળે છે, એટલે કે
$\cardfont{a} < \cardfont{b}$ જો અને માત્ર જો
$\cardfont{a} \in \cardfont{b}$. અહીં જે વસ્તુઓ
ગણાંક છે તે આપણે જાણીએ છીએ તેમના માટે ફ્રેક્ટુર
અક્ષરો વાપરીશું. આ પ્રયોગ પ્રચલિત છે. બીજી સામાન્ય
રીત ગ્રીક અક્ષરો લેવાની છે, કારણ કે ગણાંકો ક્રમવાચક
સંખ્યાઓ છે; પરંતુ એવા અક્ષરો મૂળાક્ષરના મધ્યમાંથી
પસંદ કરે છે, જેમ કે $\kappa, \lambda$.):
\begin{lem}\ollabel{lem:CardinalsBehaveRight}
કોઈ પણ ગણો $A$ અને $B$ માટે:
\begin{align*}
	\cardeq{A}{B} &\text{ iff } \card{A} = \card{B}\\
	\cardle{A}{B} &\text{ iff } \card{A} \leq \card{B}\\
	\cardless{A}{B}&\text{ iff } \card{A} < \card{B}
\end{align*}
\end{lem}

\begin{proof}
બીજા વિધાનની ડાબેથી જમણી દિશા સાબિત કરીશું
(બીજા કિસ્સા સમાન છે અને સ્વાધ્યાય માટે છોડીએ છીએ).
આ આકૃતિ જુઓ:
\begin{center}
	\begin{tikzpicture}
	\node (nodea) {$A$};
	\node[right = 6em of nodea] (nodeb) {$B$};
	\node[below = 2em of nodea] (nodecarda) {$\card{A}$};
	\node[below = 2em of nodeb] (nodecardb) {$\card{B}$};
	\draw[->] (nodea)--(nodeb);
	\draw[<->] (nodea)--(nodecarda);
	\draw[<->] (nodeb)--(nodecardb);
	\draw[->, dashed] (nodecarda)--(nodecardb);
\end{tikzpicture}
\end{center}
બંને બાજુ તીરવાળાં રેખાખંડો !!{bijection}s દર્શાવે છે;
એનું અસ્તિત્વ \olref{lem:CardinalsExist}થી મળે છે.
$\cardle{A}{B}$ ધારીએ, તો $A\to B$ એવું
!!a{injection} મળે છે. હવે તીરોને $\card{A}$થી
$A$, ત્યાંથી $B$ અને પછી $\card{B}$ સુધી અનુસરતાં
$\card{A} \to \card{B}$ એવું !!a{injection} મળે
છે (તૂટક તીર).
\end{proof}\noindent \olref{lem:CardinalsBehaveRight} વાપરીને
શ્રેડર--બર્નસ્ટાઇન પ્રમેય ફરી સાબિત કરી શકીએ.
તેનું વિધાન છે કે $\cardle{A}{B}$ અને
$\cardle{B}{A}$ હોય, તો $\cardeq{A}{B}$.
આપણે તેને \olref[sfr][siz][sb]{thm:schroder-bernstein}
તરીકે કહ્યું હતું, પરંતુ સૌપ્રથમ વધુ મહેનતથી
\olref[sfr][infinite][card-sb]{sec}માં સાબિત કર્યું હતું.
હવે જુઓ:

\begin{proof}[શ્રેડર--બર્નસ્ટાઇનની ફરી સાબિતી]
$\cardle{A}{B}$ અને $\cardle{B}{A}$ હોય, તો
\olref{lem:CardinalsBehaveRight} મુજબ
$\card{A} \leq \card{B}$ અને
$\card{B} \leq \card{A}$. તેથી ત્રિવિભાજન અને
\olref{lem:CardinalsBehaveRight}થી
$\card{A} = \card{B}$ તથા $\cardeq{A}{B}$.
\end{proof}
\noindent
આ સાબિતી બહુ સરળ છે, પરંતુ ગૂઢ રીતે તે
પ્રતિસ્થાપન પર
(\olref[ordinals][ordtype]{thmOrdinalRepresentation}
માટે) અને સુક્રમિતતા પર
(\olref{lem:CardinalsBehaveRight} માટે)
આધાર રાખે છે. તેની સામે,
\olref[sfr][infinite][card-sb]{sec}ની સાબિતી
ઘણી વધુ સ્વતંત્ર હતી (ખરેખર તે~$\Zminus$માં પણ
કરી શકાય છે).

\end{document}
